\section*{NeurIPS Paper Checklist}

\begin{enumerate}

\item {\bf Claims}
    \item[] Question: Do the main claims made in the abstract and introduction accurately reflect the paper's contributions and scope?
    \item[] Answer: \answerYes{}.
    \item[] Justification: The abstract and Section~\ref{sec:intro} explicitly state the narrow claim: the interface is operational, auditable, and active at scale, but symbolic revision is not yet statistically isolated as a within-family calibration gain.

\item {\bf Limitations}
    \item[] Question: Does the paper discuss the limitations of the work performed by the authors?
    \item[] Answer: \answerYes{}.
    \item[] Justification: Section~\ref{sec:limitations} discusses internal validation only, small TCGA sample size, hand-written rule bases, possible MC-dropout miscalibration, and the fact that this is not a full NARS implementation.

\item {\bf Theory assumptions and proofs}
    \item[] Question: For each theoretical result, does the paper provide the full set of assumptions and a complete (and correct) proof?
    \item[] Answer: \answerNA{}.
    \item[] Justification: The paper defines an interface and reports empirical results, but it does not claim new theorems. The equations and assumptions used by the method are stated in Section~\ref{sec:method}.

\item {\bf Experimental result reproducibility}
    \item[] Question: Does the paper fully disclose all the information needed to reproduce the main experimental results of the paper to the extent that it affects the main claims and/or conclusions of the paper (regardless of whether the code and data are provided or not)?
    \item[] Answer: \answerYes{}.
    \item[] Justification: Section~\ref{sec:experiments} reports the datasets, split sizes, model family, baselines, MC-dropout passes, and metrics. The anonymous repository is linked in the data and code statement.

\item {\bf Open access to data and code}
    \item[] Question: Does the paper provide open access to the data and code, with sufficient instructions to faithfully reproduce the main experimental results, as described in supplemental material?
    \item[] Answer: \answerYes{}.
    \item[] Justification: The anonymous repository is available at \url{https://anonymous.4open.science/r/NGTA/}. The benchmarks are public datasets, and the repository contains reproduction materials for review.

\item {\bf Experimental setting/details}
    \item[] Question: Does the paper specify all the training and test details (e.g., data splits, hyperparameters, how they were chosen, type of optimizer) necessary to understand the results?
    \item[] Answer: \answerYes{}.
    \item[] Justification: Section~\ref{sec:experiments} reports the main split sizes, input dimensions, transformer architecture, uncertainty configuration, baselines, and evaluation metrics. Lower-level run details are provided in the anonymous repository.

\item {\bf Experiment statistical significance}
    \item[] Question: Does the paper report error bars suitably and correctly defined or other appropriate information about the statistical significance of the experiments?
    \item[] Answer: \answerYes{}.
    \item[] Justification: Tables~\ref{tab:tcga_results} and \ref{tab:wids_results} report bootstrap confidence intervals for AUC, Brier score, and ECE. Section~\ref{sec:results} states when paired bootstrap intervals include zero.

\item {\bf Experiments compute resources}
    \item[] Question: For each experiment, does the paper provide sufficient information on the computer resources (type of compute workers, memory, time of execution) needed to reproduce the experiments?
    \item[] Answer: \answerNo{}.
    \item[] Justification: The anonymous repository provides the executable reproduction code and run entry points, but the paper does not list exact hardware, memory, or wall-clock runtime. The experiments are small tabular benchmark runs, and exact compute accounting can be added if the selected workshop explicitly requires it.

\item {\bf Code of ethics}
    \item[] Question: Does the research conducted in the paper conform, in every respect, with the NeurIPS Code of Ethics \url{https://neurips.cc/public/EthicsGuidelines}?
    \item[] Answer: \answerYes{}.
    \item[] Justification: The work uses public, de-identified benchmark data, makes no deployment claim, and frames the system as requiring external validation and clinical review before any clinical use.

\item {\bf Broader impacts}
    \item[] Question: Does the paper discuss both potential positive societal impacts and negative societal impacts of the work performed?
    \item[] Answer: \answerYes{}.
    \item[] Justification: The paper discusses the positive goal of auditable uncertainty-aware clinical prediction and the negative risks of miscalibrated uncertainty, weak rules, and insufficient validation in Section~\ref{sec:limitations}.

\item {\bf Safeguards}
    \item[] Question: Does the paper describe safeguards that have been put in place for responsible release of data or models that have a high risk for misuse (e.g., pre-trained language models, image generators, or scraped datasets)?
    \item[] Answer: \answerNA{}.
    \item[] Justification: The paper does not release high-risk generative models, scraped datasets, or deployment-ready medical devices. The released research code is tied to public, de-identified benchmarks.

\item {\bf Licenses for existing assets}
    \item[] Question: Are the creators or original owners of assets (e.g., code, data, models), used in the paper, properly credited and are the license and terms of use explicitly mentioned and properly respected?
    \item[] Answer: \answerNo{}.
    \item[] Justification: The paper identifies the public benchmark sources and cites related dataset/model literature, and the repository documents expected data files and preprocessing. However, the project currently has no explicit license file, so code/license terms should be finalized before public release or camera-ready submission.

\item {\bf New assets}
    \item[] Question: Are new assets introduced in the paper well documented and is the documentation provided alongside the assets?
    \item[] Answer: \answerYes{}.
    \item[] Justification: The main new asset is research code for the dynamic evidential routing pipeline. The anonymous repository provides documentation and reproduction materials for review.

\item {\bf Crowdsourcing and research with human subjects}
    \item[] Question: For crowdsourcing experiments and research with human subjects, does the paper include the full text of instructions given to participants and screenshots, if applicable, as well as details about compensation (if any)?
    \item[] Answer: \answerNA{}.
    \item[] Justification: The paper does not involve crowdsourcing or new human-subjects experiments.

\item {\bf Institutional review board (IRB) approvals or equivalent for research with human subjects}
    \item[] Question: Does the paper describe potential risks incurred by study participants, whether such risks were disclosed to the subjects, and whether Institutional Review Board (IRB) approvals (or an equivalent approval/review based on the requirements of your country or institution) were obtained?
    \item[] Answer: \answerNA{}.
    \item[] Justification: The work uses public, de-identified benchmark data and does not collect new participant data. Therefore IRB approval is not applicable for the reported experiments.

\item {\bf Declaration of LLM usage}
    \item[] Question: Does the paper describe the usage of LLMs if it is an important, original, or non-standard component of the core methods in this research? Note that if the LLM is used only for writing, editing, or formatting purposes and does \emph{not} impact the core methodology, scientific rigor, or originality of the research, declaration is not required.
    \item[] Answer: \answerYes{}.
    \item[] Justification: The ethics, data, and code availability statement discloses that LLMs were used to assist with drafting, mathematical formulation, and code generation. The author(s) independently checked the derivations, implementation, and reported results and take full responsibility for the work.

\end{enumerate}
